% ─────────────────────────────────────────────────────────────────
%  Draft of §IV.A — to replace existing §IV.A (Phase diagram, lines
%  362–429) and §IV.B (Spatiotemporal structure, lines 431–485).
%  Hot-runaway attractor is described later in §IV.D, not here.
% ─────────────────────────────────────────────────────────────────

\section{Results and Discussion}
\label{sec:results}

The dispersion analysis of Sec.~\ref{sec:lsa} establishes the linear
bifurcation conditions for the homogeneous swollen state.  We now
turn to direct simulations of the full PDE system
\eqref{eq:governing} in the planar slab and report four results.
First, default-IC simulations across the principal parameter planes
realize three distinct attractors—cold-swollen, an LCST-front
oscillation, and a frozen partial-collapse front
(Sec.~\ref{sec:phasediagram}).  Second, the LCST-front attractor is
a relaxation oscillator on a slow manifold generated by the surface
boundary condition together with the $\varphi$-dependent diffusion
barrier; the same barrier prevents the gel from undergoing global
collapse (Sec.~\ref{sec:lcst_front_mechanism}).  Third, the 0D
linear stability analysis is a necessary but not sufficient
predictor of the PDE oscillation: a closed-form parameter
inversion of the homogeneous-SS conditions identifies a Whitney-
cusp catastrophe whose collapsed-only branch covers
$\sim\!51\%$ of the displayed $(\Bi_T,S_\chi)$ plane, and the
saddle--node-on-invariant-circle scaling at the upper-$S_\chi$
boundary confirms that the cycle is destroyed by a global
topological event rather than by a local Hopf instability of any
homogeneous branch (Sec.~\ref{sec:lsa_vs_nl}).  Fourth,
initial-condition sweeps reveal a fourth attractor—a hot-runaway
super-swollen steady state—coexisting with the LCST-front cycle in
a substantial portion of the C-region; the bistability gives a
falsifiable IC-induced hysteresis prediction
(Sec.~\ref{sec:multistability}).

\subsection{Realized attractors and their spatiotemporal structure}
\label{sec:phasediagram}

We integrate the PDE system from the canonical cold-swollen initial
condition ($J_0=1.30$, $\theta_0=0$, $u_0=0.02$) to dimensionless
time $t=200$, classify each parameter point by combining the surface
swelling amplitude, the peak polymer fraction, and the oscillation
period into a six-class regime label, and sweep $25\times25$ points
in the $(\Bi_T, S_\chi)$ plane and $25\times15$ points in the
$(\Bi_T, \Da)$ plane (Fig.~\ref{fig:fig4_phase}).

\paragraph{Realized attractors (default-IC).}
Three of the six classification labels are populated:
\begin{itemize}
\item[(i)] a uniform \emph{cold-swollen} steady state at the cold
  edge of the diagram (low $S_\chi$ or high $\Bi_T$);
\item[(ii)] a large-amplitude \emph{LCST-front} oscillation occupying
  the central $\sim\!45\%$ of the $(\Bi_T,S_\chi)$ plane, in which a
  narrow surface shell cycles between collapsed and swollen states
  while the centre stays cold-swollen;
\item[(iii)] a \emph{frozen partial-collapse} state at large $\Bi_T$,
  with the same two-zone spatial structure as the LCST-front cycle
  but stationary in time.
\end{itemize}
The remaining classifier-allowed regimes (globally collapsed steady,
globally collapsing oscillation) never appear in the $1{,}000$
default-IC simulations covered by the two parameter sweeps; the
mechanism that suppresses them is discussed in
Sec.~\ref{sec:lcst_front_mechanism}.  A fourth attractor, a uniform
\emph{hot-swollen runaway} steady state, exists in a substantial
sub-region of the LCST-front map but is inaccessible from the
default cold IC and is revealed only by initial-condition variation
(Sec.~\ref{sec:multistability}).

\paragraph{Spatiotemporal structure of the three default-IC
attractors.}
Figure~\ref{fig:fig_all_regimes} shows the spatiotemporal fields
$J(\xi,t)$, $\theta(\xi,t)$ and the surface and centre $J(t)$ traces
for one representative point per attractor.

The cold-swollen state [Fig.~\ref{fig:fig_all_regimes}(a)] is
spatially uniform with $J\!\simeq\!1.30$ and $\theta\!\simeq\!0$
everywhere; the surface BC is passively satisfied without driving
significant transport.

The LCST-front cycle [Fig.~\ref{fig:fig_all_regimes}(b)] has a
two-zone structure.  The outer shell ($\xi\!\in\![0.885,\,1]$ at
the working point) cycles between the swollen state
($J\!\sim\!1.4$) and the collapsed state ($J\!\sim\!0.15$) with
period $T\!\simeq\!8\text{--}11$.  The inner core ($\xi<0.885$)
stays at $J\!\simeq\!1.39$ throughout the cycle.  The temperature
field tracks the asymmetry: $\theta_\mathrm{surf}$ swings between
$\sim\!1.5$ and $\sim\!4.8$ each cycle, while
$\theta_\mathrm{centre}$ is held nearly constant at
$\theta_\mathrm{centre}\!\simeq\!3.2$ by heat conduction from the
warm shell.

The frozen-front state [Fig.~\ref{fig:fig_all_regimes}(c)] inherits
the same two-zone spatial structure, now stationary in time
($J_\mathrm{surf}\!\simeq\!0.15$,
$J_\mathrm{centre}\!\simeq\!1.4$).  At sufficiently large $\Bi_T$,
the surface heat flux $\Bi_T\theta_\mathrm{surf}$ is large enough
to clamp $\theta_\mathrm{surf}$ below the ignition threshold and
the cycle of Sec.~\ref{sec:lcst_front_mechanism} is locked at
phase~3.
% ─────────────────────────────────────────────────────────────────
%  Draft of §IV.B (relaxation oscillator) — to replace existing
%  §IV.C (Mechanism of the LCST-front oscillation, lines 487–555 of
%  main.tex).  Uses figures B2 (slow_manifold), B3 (period_scaling),
%  and the existing time-series / heatmap panels of current Fig.4.
% ─────────────────────────────────────────────────────────────────

\subsection{A relaxation oscillator on a spatially-generated slow
manifold}
\label{sec:lcst_front_mechanism}

The LCST-front cycle is a relaxation oscillator with two distinguishing
features: (i)~the fast variable is mechanical ($J$), the slow
variables are thermal and chemical ($\theta$, $u$); and (ii)~the
slow manifold is generated by the spatial structure of the gel rather
than by an autonomous 0D nonlinearity.  Each cycle traverses two
stable branches of this manifold connected by fast jumps, and the
period is dominated by the slowest of the four phases.

\paragraph{Fast--slow decomposition.}  At the working point the
mechanical relaxation timescale set by the dimensionless mobility
$\Ms_0$ is roughly two orders of magnitude shorter than the thermal
($1/\Bi_T$) and reactant-supply ($1/\Bi_c$) timescales.  Numerically,
the speed $|d(\theta_\mathrm{surf}, J_\mathrm{surf})/dt|$ along the
limit cycle [Fig.~\ref{fig:fig2_mechanism}(c)] varies by a factor of
$\sim\!600$ between the slow drifts and the fast jumps.  We may
therefore treat $J$ as a fast variable and follow the standard
relaxation-oscillator construction.

\paragraph{Spatially generated slow manifold.}
Quasi-static balance of the surface mechanical equation with the
Robin boundary condition $\Bi_\mu(\mu - \mu_b) = 0$ gives the locus
\begin{equation}
  \mu(J,\theta) \;=\; \mu_b
  \label{eq:slow_manifold}
\end{equation}
in the $(\theta, J)$ plane.  The Flory--Huggins--Rehner free energy
is non-monotonic in $J$ for sufficiently large $\chi_\mathrm{eff}=
\chi_\infty + S_\chi\theta + \chi_1 \phi$ (above the LCST), so this
locus has the canonical S-shape
[Fig.~\ref{fig:fig2_mechanism}(c)]: an upper \emph{swollen} branch
($J \gtrsim 0.8$) for $\theta \le \theta_\mathrm{up}$, a lower
\emph{collapsed} branch ($J \lesssim 0.2$) for
$\theta \ge \theta_\mathrm{lo}$, and a saddle (unstable middle
branch) connecting the two folds. At the working point,
$\theta_\mathrm{lo}\!\simeq\!0.85$ and
$\theta_\mathrm{up}\!\simeq\!3.18$, defining a bistable region
$\theta\in[\theta_\mathrm{lo},\theta_\mathrm{up}]$ in surface
temperature.  Crucially, this slow manifold lives at the surface and
is enforced by the surface BC; it has no analog in a 0D
homogeneous-state reduction, where $\mu = \mu_b$ would have to be
imposed everywhere simultaneously together with the heat and reactant
balances (Sec.~\ref{sec:lsa_vs_nl}).

\paragraph{Cycle traversal.}  The PDE limit cycle, projected onto
$(\theta_\mathrm{surf}, J_\mathrm{surf})$, traces the slow manifold's
two stable branches connected by near-vertical fast jumps
[Fig.~\ref{fig:fig2_mechanism}(c)].  We label four canonical phases:

\begin{enumerate}[leftmargin=*,label=\arabic*.]
\item \textbf{Ignition (slow, on swollen branch).}  The surface starts
  at the cold end of the swollen branch ($\theta\!\simeq\!1.5$,
  $J\!\simeq\!1.4$) with abundant reactant
  ($u_\mathrm{surf}\!\simeq\!1$).  A small thermal fluctuation
  amplifies the Arrhenius factor; the heat source $\Da\,J\hat R$
  exceeds the surface heat loss $\Bi_T\theta_\mathrm{surf}$ and
  $\theta_\mathrm{surf}$ rises monotonically.
\item \textbf{Collapse (fast).}  As $\theta_\mathrm{surf}$ exceeds the
  upper-fold value, the swollen branch ceases to exist; $J$ falls
  rapidly along the fast manifold from $J\!\sim\!0.8$ to
  $J\!\sim\!0.15$ over a small fraction of the cycle period.
\item \textbf{Quenching (slow, on collapsed branch).}  With
  $\varphi\!\sim\!1$ in the collapsed shell the diffusion factor
  $D(\varphi) \propto (1-\varphi)^2$ drops by two orders of magnitude,
  throttling reactant supply from the bath.  Boundary-layer reactant
  is consumed by the residual reaction; with the heat source gone,
  $\theta_\mathrm{surf}$ relaxes towards the bath temperature on the
  surface heat-loss timescale $1/\Bi_T$.  This is generally the
  longest phase of the cycle.
\item \textbf{Re-swelling (fast).}  When $\theta_\mathrm{surf}$ falls
  below the lower-fold value, the collapsed branch terminates and $J$
  jumps back up to the swollen branch.  The diffusion barrier
  reopens, the bath rapidly resupplies reactant, and the cycle
  returns to phase~1.
\end{enumerate}

The slow drifts (phases 1 and 3) and the fast jumps (phases 2 and 4)
are clearly resolved both as time-series plateaus and ramps
[Fig.~\ref{fig:fig2_mechanism}(b)] and as differently-shaded
segments of the limit cycle in the slow-manifold plane
[Fig.~\ref{fig:fig2_mechanism}(c)].

\paragraph{Self-limiting collapse.}  The same $\varphi$-dependent
diffusion barrier that quenches the reaction at the surface during
phase~3 also prevents collapse from invading the gel interior.  The
LCST front penetrates only to
$\xi_\mathrm{LCST}\!\simeq\!0.885$ at the working point; the inner
$88.5\%$ of the slab never crosses the LCST and remains swollen
throughout the cycle.  This barrier explains why the classifier-
allowed regimes \emph{globally collapsed steady} and \emph{globally
collapsing oscillation} are absent from $1{,}000$ default-IC
simulations spanning the parameter space (Sec.~\ref{sec:phasediagram}):
once any portion of the gel collapses, the inner region is shielded
from further reactant supply and cannot ignite. The same mechanism
is, in this sense, a self-limiting protective feature of this class
of thermo-LCST gel against runaway global collapse despite the
strongly exothermic Arrhenius
kinetics~\cite{FrankKamenetskii1969,Uppal1974}.

\paragraph{Period structure.}  Across the LCST-front regime
[Fig.~\ref{fig:fig2_mechanism}(e)], the cycle period $T_\mathrm{PDE}$
spans roughly $6$--$15$, with a minimum of $\sim\!8$ near the
``sweet-spot'' $S_\chi\!\simeq\!0.75$ and divergence at the upper-
$S_\chi$ boundary where the cycle disappears through a saddle--node-
on-invariant-circle (SNIC)-like bifurcation
(Sec.~\ref{sec:lsa_vs_nl}).  The cooling timescale $1/\Bi_T$ sets the
order of magnitude of the period (prefactor $\mathcal{O}(1)$
absorbing the ignition delay and the spatial barrier-release time),
but the period is not a simple power law of $\Bi_T$ alone because
phase~3 also depends on the residual interior heat conduction and on
the timescale at which the boundary BC restores
$u_\mathrm{surf}$ once the barrier reopens.
% ─────────────────────────────────────────────────────────────────
%  Draft of §IV.C (Linear stability is necessary but not sufficient)
%  — slim version after retiring C1(b) (PDE-overlay panel) and C2
%  (T_PDE / T_LSA contrast).  Argument now leans on §IV.B's
%  spatially-generated slow manifold; C1(a) (analytic SS-existence
%  map with Whitney-cusp fold) and C3 (SNIC period-divergence scaling)
%  are the two empirical anchors.
% ─────────────────────────────────────────────────────────────────

\subsection{Linear stability is necessary but not sufficient}
\label{sec:lsa_vs_nl}

Section~\ref{sec:lcst_front_mechanism} showed that the LCST-front
attractor is a relaxation cycle whose slow manifold is generated by
the surface boundary condition together with the
$\varphi$-dependent diffusion barrier.  No spatially uniform reduction
can produce that manifold: at the surface the locus
$\mu(J,\theta)=\mu_b$ is enforced by the Robin BC, while in the
interior the same locus is irrelevant because the bulk equilibrium is
buffered by transport.  Linear stability of an available homogeneous
branch is therefore at best a \emph{necessary} condition for the PDE
to oscillate; it cannot be sufficient because the cycle traverses
two branches of a slow manifold that does not exist in the 0D
reduction.  This subsection makes the necessary-condition statement
quantitative on two complementary axes: (i)~where the homogeneous SS
even exists as a real root, and (ii)~how the cycle is destroyed at
the parameter-plane boundaries of the LCST-front regime.

\paragraph{Existence of the homogeneous swollen branch (analytic).}
The 0D steady-state system \eqref{eq:lsa_ss} admits an explicit
parameterization: at any pair $(J,\theta)$ with $J>\varphi_{p0}$,
the two SS conditions solve in closed form for the dependent
control parameters,
\begin{equation}
  S_\chi(J,\theta) =
    \frac{\mu_b - \mu_{\mathrm{base}}(J,\theta)}{\theta\,\varphi^2}, \qquad
  \Bi_T(J,\theta) = \frac{K(J,\theta)}{\theta\,(1+K/\Bi_c)},
  \label{eq:ss_param_inv}
\end{equation}
where $\mu_{\mathrm{base}}$ is the part of $\mu$ that excludes the
$S_\chi\theta\,\varphi^2$ coupling and
$K(J,\theta)=\Da\,J\,(1-\varphi)^{m_\mathrm{act}}\,T(\theta)$.  The
entire 0D SS surface is therefore the image of the $(J,\theta)$ plane
under \eqref{eq:ss_param_inv}, and the saddle--node fold locus in
$(\Bi_T, S_\chi)$ is exactly the image of the curve
$\det\partial(F_1,F_2)/\partial(J,\theta)=0$.  Numerically that locus
is a single connected component in $(J,\theta)$; under the map
\eqref{eq:ss_param_inv} it self-folds in $(\Bi_T,S_\chi)$ into two
branches meeting at a Whitney cusp at
$(\Bi_T,S_\chi)\!\simeq\!(0.43, 2.17)$ — a cusp catastrophe in the
homogeneous-SS surface.

Counting the real roots of \eqref{eq:lsa_ss} cell-by-cell on a
$360\times 240$ grid in $(\Bi_T, S_\chi)$ via the closed-form
inverse \eqref{eq:ss_param_inv} (one sign-change scan per cell along
$J$, exact and ${\sim}10^{2}\times$ faster than fsolve) gives
Fig.~\ref{fig:fig_homog_ss}.  The fraction of the displayed plane
occupied by each multiplicity is summarised below
\begin{center}
\begin{tabular}{lr}
\hline
Region & Area fraction \\
\hline
swollen branch only & $37\%$ \\
collapsed branch only & $51\%$ \\
bistable (3 roots: 2 stable + 1 saddle) & $11\%$ \\
no root (off the displayed window) & $0.3\%$ \\
\hline
\end{tabular}
\end{center}
exactly $0\%$ of cells have $2$ roots: this odd-only multiplicity is
forced by the fold topology (any horizontal line in
$(\Bi_T, S_\chi)$ at fixed $S_\chi$ crosses the two fold arcs an even
number of times, giving root counts $1, 3, 5, \ldots$).  The default
working point $(\Bi_T=0.10,S_\chi=1.0)$ and the canonical C-cell
$(\Bi_T=0.059, S_\chi=1.80)$, both used elsewhere in this section,
sit deep inside the ``collapsed-only'' region: at neither point does
the swollen homogeneous SS that one would linearise about exist as a
real root, so a 0D LSA about the swollen branch is structurally
unavailable.  The PDE nevertheless finds the LCST-front cycle there.
The mismatch is the
necessary-but-not-sufficient statement made geometric: the
spatially-generated slow manifold of
Sec.~\ref{sec:lcst_front_mechanism} continues to exist past the
saddle--node fold of the homogeneous problem because its swollen
branch is enforced \emph{at the surface} by the Robin BC, not
required to coexist with bulk reactant balance.

\paragraph{SNIC scaling at the upper-$S_\chi$ exit.}
At the upper-$S_\chi$ boundary of the LCST-front regime
[Fig.~\ref{fig:fig4_phase}(e)], the cycle period diverges as
$S_\chi\to S_\chi^{(c)}$ from below.  We test the period scaling
along a 1D slice through this boundary against the two canonical
infinite-period scenarios
\begin{equation}
  T_\mathrm{SNIC}(S_\chi)\propto (S_\chi^{(c)}-S_\chi)^{-1/2},
  \qquad
  T_\mathrm{hom}(S_\chi)\propto -\ln(S_\chi^{(c)}-S_\chi),
  \label{eq:snic_vs_hom}
\end{equation}
using the same 1D scan
[Fig.~\ref{fig:fig_snic_scaling}].  A two-parameter least-squares fit
of $T_\mathrm{PDE}(S_\chi)$ on the bracket
$S_\chi\in[1.10, 1.32]$ yields a SNIC residual sum of squares
$\mathrm{RSS}_\mathrm{SNIC}=0.068$ versus
$\mathrm{RSS}_\mathrm{hom}=0.152$, with the SNIC fit giving the
critical value $S_\chi^{(c)}\!\approx\!1.34$.  The SNIC inverse-square-
root law captures the data more than twice as well as the homoclinic
log law, consistent with the cycle being destroyed by the collision
of a saddle and a node on the limit cycle rather than by a homoclinic
orbit hitting a saddle off the cycle.  In either reading the cycle's
disappearance is a global topological event in the surface
$(\theta_\mathrm{surf}, J_\mathrm{surf})$ slow-manifold plane and not
a local linear bifurcation of any single homogeneous branch — the
expected signature of a relaxation oscillator on a spatially
generated slow manifold.

\paragraph{Summary.}
The 0D linear stability analysis carries a precise but limited
informational content: it identifies the local stability of available
homogeneous branches, but the swollen homogeneous branch fails to
exist over $\sim\!51\%$ of the parameter plane covered here, and even
where it does the cycle that the PDE realises is generically not the
small-amplitude continuation of any homogeneous Hopf bifurcation.
The dominant nonlinear behaviour of this thermo-LCST gel is a
relaxation oscillator on a spatially generated slow manifold, born
and destroyed by global topological events
(Sec.~\ref{sec:lcst_front_mechanism} and Eq.~\ref{eq:snic_vs_hom}).
Whether the same parameter plane harbours a hidden second attractor
not visible to either linear analysis or default-IC numerical scans
is settled in Sec.~\ref{sec:multistability}.
% ─────────────────────────────────────────────────────────────────
%  Draft of §IV.D (hidden hot-runaway + IC hysteresis prediction) —
%  to replace existing §IV.E (Hidden multistability, lines 609–696
%  of main.tex).  Pulls in figure D2 (hysteresis_C_theta0) and
%  D1 (refined basin at C-cell), and inherits Fig.6 (a),(b),(c).
% ─────────────────────────────────────────────────────────────────

\subsection{A hidden hot-runaway state and the IC-induced hysteresis
prediction}
\label{sec:multistability}

Sec.~\ref{sec:lsa_vs_nl} closed by asking whether the C-region of
the parameter plane—where the swollen homogeneous SS does not exist
yet the PDE oscillates—corresponds to genuine attractor multiplicity
of the spatially extended problem.  We resolve this by initial-
condition (IC) sweeps at a representative C-cell
($\Bi_T=0.059$, $S_\chi=1.80$).

\paragraph{Two co-existing attractors.}
At identical bath conditions and reaction parameters, simulations
initialized in different regions of $(J_0, \theta_0)$ space converge
to one of two qualitatively distinct attractors
[Fig.~\ref{fig:fig_manifold}(c)]: (i)~the LCST-front cycle of
Sec.~\ref{sec:lcst_front_mechanism}, with bulk-mean
$\langle J\rangle\!\approx\!2$,
$\langle\theta\rangle\!\approx\!2$, and (ii)~a previously-unreported
\emph{hot-runaway} steady state with
$\langle J\rangle\!\approx\!5$,
$\langle\theta\rangle\!\approx\!12$.
The two attractors differ by
$\Delta\langle J\rangle\!\approx\!3.3$ and
$\Delta\langle\theta\rangle\!\approx\!10$, far beyond any numerical
tolerance, so the bistability is genuine.

\paragraph{Spatial structure and mechanism of the hot-runaway state.}
The hot-runaway state [Fig.~\ref{fig:fig_manifold}(b)] is spatially
uniform with $J\!\simeq\!5.2$, $\theta\!\simeq\!12$,
$\varphi\!\simeq\!0.026$, and a thin reaction layer at the surface
where reactant is consumed almost entirely
($u_\mathrm{surf}\!\sim\!10^{-7}$).  The mechanism is diffusion-
limited combustion: the high-$\theta$ Arrhenius factor exceeds
$10^5$ and burns through any reactant in the boundary layer, but the
surface BC continues to deliver reactant at the diffusion-limited
rate $\Bi_c$.  The gel remains super-swollen because the high-$J$
Flory--Huggins minimum coexists with the collapsed minimum whenever
$\chi_\mathrm{eff}$ is sufficiently large; at $\varphi\!\sim\!0.026$
the $\chi_\mathrm{eff}\,\varphi^2$ penalty is negligible and the
elastic term $\Omega_e\,(J-1/J)$ becomes the dominant contribution to
the chemical potential.  The choice between the two
Flory--Huggins minima is set by the IC: trajectories landing in the
basin of the swollen branch reach hot-runaway, while trajectories
landing in the basin of the collapsed branch are ejected to the
LCST-front cycle by the surface BC and the diffusion-barrier
feedback.

\paragraph{Basin asymmetry and why the default sweep misses the
runaway.}
A refined IC sweep on a $15\times 15$ grid at the C-cell
[Fig.~\ref{fig:fig_manifold}(c)] places the basin boundary at
roughly $J_0\!\approx\!1.0$, $\theta_0\!\approx\!1$: only ICs in the
upper-left quadrant of $(J_0, \theta_0)$ space (cold, swollen) reach
the cycle.  ICs with any of (i)~$J_0\!\lesssim\!0.4$ (pre-collapsed),
(ii)~$\theta_0\!\gtrsim\!3$ (ignited), or (iii)~combinations
spanning the basin boundary fall into the runaway.  The basin of
the cycle covers
$\sim\!13\%$ of the IC grid (16 of the 119 successfully integrated
cells of an $11{\times}11$ grid spanning $J_0\in[0.20,1.30]$,
$\theta_0\in[0,10]$), a small fraction of the total. This asymmetry explains why the default
cold-IC sweep of Fig.~\ref{fig:fig4_phase} sees only the cycle and
never the runaway: the default cold IC sits well inside the cycle's
basin at every grid point of the parameter sweep.

\paragraph{IC-induced hysteresis: a falsifiable prediction.}
A 1D slice through the basin space gives a quantitative,
experimentally accessible signature of the bistability
[Fig.~\ref{fig:fig_manifold}(d)].  Holding all reaction and bath
parameters fixed at the C-cell value with the initial reactant
loading $u_0=0.5$, we vary the initial thermal pulse $\theta_0$ from
$0$ to $6$ for two seed conditions:
\begin{itemize}
\item ``cold'' seed ($J_0=1.30$): initially near the cold-swollen
  homogeneous state.  For $\theta_0$ below the threshold
  $\theta_0^{\!*}\!\approx\!2.4$ the trajectory remains in the
  cycle's basin and
  $\langle J\rangle_\mathrm{late}\!\approx\!2.0$;
  for $\theta_0$ above the threshold the seed crosses the basin
  boundary and
  $\langle J\rangle_\mathrm{late}\!\approx\!5.9$,
  indistinguishable from the runaway state.  The transition is
  sharp—the discontinuity in $\langle J\rangle_\mathrm{late}$ falls
  within a single $\Delta\theta_0=0.25$ sampling step.
\item ``pre-collapsed'' seed ($J_0=0.20$): initialized inside the
  runaway basin; for every $\theta_0$ in the scanned range the
  trajectory terminates in the runaway-type state with
  $\langle J\rangle_\mathrm{late}\!\approx\!4.2\text{--}5.5$,
  independent of $\theta_0$.
\end{itemize}
The split between the two terminal-state curves over the bistable
$\theta_0$ interval $[0, 2.4]$ is the hysteresis signature.
Experimentally, this would manifest as a sharp jump in long-time
gel volume (or temperature) when the initial perturbation magnitude
is varied across $\theta_0^{\!*}$, with no return jump on the
reverse sweep so long as the C-cell parameters are held fixed.
This is a quantitative prediction of the model that the linear
theory does not provide and that distinguishes spatially-extended
thermo-LCST gels from the well-mixed CSTR
analog~\cite{Uppal1974}.

\paragraph{Summary.}
The PDE attractor diagram is therefore richer than the single-IC
sweep of Fig.~\ref{fig:fig4_phase} suggests, and richer than 0D LSA
can predict.  The hot-runaway state is invisible to both linear
analysis and to default-IC numerical scans, and its appearance is
contingent on the structure of the basin of attraction in IC space.
Combined with the slow-manifold relaxation cycle of
Sec.~\ref{sec:lcst_front_mechanism}, this completes a four-attractor
inventory of the spatially extended problem (cold-swollen, LCST-
front cycle, frozen front, hot-runaway), whose existence and
stability are governed not by any single homogeneous-state
linearization but by the spatial coupling between the Flory--Huggins
free energy and the $\varphi$-dependent diffusion barrier.
